% &%%%%%%%%%%%%%%%%%%%%%%%%%%%%%%%%%%%%%%%%%%%%%%%%%%%%%%%%%%%
% &%%                                                       %%
% &%%                     NOMENCLATURE                      %%
% &%%                                                       %%
% &%%%%%%%%%%%%%%%%%%%%%%%%%%%%%%%%%%%%%%%%%%%%%%%%%%%%%%%%%%%

% !===========================================================

\chapter*{Nomenclature}

\mtcaddchapter%                                                 % Ajout du chapitre en tant que chapitre fantôme (pour ne pas lui attribuer de numéro)
\addcontentsline{toc}{chapter}{Nomenclature}                    % Ajout du chapitre à la ToC
\markboth{Nomenclature}{Nomenclature}                           % Permet d'afficher le chapitre en en-tête

% !===========================================================

\textbf{Sauf mention contraire, les notations utilisées dans ce document sont celles définies ci-dessous.}

\glsfindwidesttoplevelname%                                     % Calcule la largeur du terme le plus long
\setglossarystyle{alttree}                                      % Définit le style du glossaire

% Vous pouvez rajouter des catégories de glossaire ici

{
\parskip=0pt

% ?-----------------------------------------------------------
% ? ACRONYMES
% ?-----------------------------------------------------------
\renewcommand*{\glossarysection}[2][]{\section*{Acronymes}}
\printglossary[type=\acronymtype]

% ?-----------------------------------------------------------
% ? LETTRES ROMANES
% ?-----------------------------------------------------------
\renewcommand*{\glossarysection}[2][]{\section*{Lettres romanes}}
\printglossary[type=roman]

% ?-----------------------------------------------------------
% ? LETTRES GRECQUES
% ?-----------------------------------------------------------
\renewcommand*{\glossarysection}[2][]{\section*{Lettres grecques}}
\printglossary[type=grec]

% ?-----------------------------------------------------------
% ? OPERATEURS ET SYMBOLES
% ?-----------------------------------------------------------
\renewcommand*{\glossarysection}[2][]{\section*{Opérateurs et Fonctions}}
\printglossary[type=operator]

}
